\clearpage
\section{Universal Hecke composition and physical exchange}
\label{sec:f1-limit-composition}

The positive Hecke tower contains more than its familiar partial-flag
corners.  Closing the tower under finite sums and self-adjoint retracts gives
a finite graded composition category, while passing to finite-dimensional
right modules gives the same category algebraically.  The tensor product is
induction along the ordered block inclusion.  This point matters: the plain
vector-space tensor product of modules over two different Hecke algebras is
not yet a module over the next Hecke algebra.

There are also two different notions of exchange.  The usual Hecke braid is
an algebraic isomorphism, but it is not unitary for the coefficient-trace
$C^*$-structure away from $q=1$.  Polar decomposition of the longest Hecke
elements instead produces canonical unitary reversals.  Their block
commutors obey cactus, rather than braid, coherence.

\subsection{The finite completion of a multiplicative sequence}

\begin{definition}[Multiplicative sequence, skeleton, and variance]
A multiplicative sequence over $\mathbb C$ is a family of unital algebras
$A_*=(A_n)_{n\geq0}$, with $A_0=\mathbb C$, and unital algebra maps
\[
 \mu_{m,n}:A_m\otimes A_n\longrightarrow A_{m+n}
\]
whose two composites from $A_l\otimes A_m\otimes A_n$ to $A_{l+m+n}$
agree.  The unit maps are the scalar identifications
$\mu_{0,n}(\lambda\otimes a)=\lambda a$ and
$\mu_{n,0}(a\otimes\lambda)=\lambda a$.

Its skeleton $\mathbb C[A_*]$ has objects $[n]$, zero morphisms between
unequal degrees, $\operatorname{End}([n])=A_n$, composition given by
algebra multiplication, and tensor induced by $\mu$.  A linear presheaf is
a functor $\mathbb C[A_*]^{\mathrm{op}}\to\mathrm{Vect}_{\mathbb C}$.
At degree $n$ it is a right $A_n$-module:
\[
 m\cdot a=F(a)m.
\]
The representable $\operatorname{Hom}(-,[n])$ is the right regular module,
because precomposition sends $x$ to $xa$.
\end{definition}
\begin{scope}
The direction of the module action is part of the definition.  No opposite
algebra is inserted later, and no morphisms between different degrees are
created at the skeletal stage.
\end{scope}
\provenance{D1221}{theory/sidequests/f1-limit/universal.md}{structural proof; no dedicated finite gate}{theory/verdicts/f1-limit-adjudication.md}

\begin{definition}[Finite Day--Schur--Weyl completion]
Assume every $A_n$ is finite-dimensional over $\mathbb C$.  Define
\[
 \operatorname{SW}_{\mathrm{fin}}(A_*)
   =\bigoplus_{n\geq0}\operatorname{mod}_{\mathrm{fd}}\!\text{-}A_n,
\]
where objects have finite degree support.  For a right $A_m$-module $M$ and
a right $A_n$-module $N$, set
\[
 M\star N=(M\otimes_{\mathbb C}N)
   \otimes_{A_m\otimes A_n}A_{m+n}.
\]
Here $A_{m+n}$ is a left $A_m\otimes A_n$-module through $\mu_{m,n}$ and a
right module by multiplication.  The unit is $\mathbb C=A_0$, and the
associator is the canonical balanced-tensor-product isomorphism.
\end{definition}
\begin{scope}
Finite-dimensionality of every level algebra is required for this finite
subcategory to be closed under $\star$.  The tensor is induction, not the
unextended tensor product $M\otimes_{\mathbb C}N$.
\end{scope}
\provenance{D1222}{theory/sidequests/f1-limit/universal.md}{structural balanced-tensor proof}{theory/verdicts/f1-limit-adjudication.md}

\begin{definition}[Finite projective composition category]
The additive Karoubi completion $\operatorname{Proj}(A_*)$ has degree-$n$
objects $(r,e)$, where $r\geq0$ and $e\in M_r(A_n)$ is an idempotent.  It
represents the right module $eA_n^r$.  For two objects in degree $n$,
\[
 \operatorname{Hom}((r,e),(s,f))=fM_{s,r}(A_n)e,
\]
with matrix multiplication as composition; unequal degrees have zero Hom.
Tensor uses the matrix amplification of $\mu_{m,n}$ and the idempotent
$\mu_{m,n}^{(r,s)}(e\otimes f)$.

If every $A_n$ is a finite-dimensional $C^*$-algebra and every $\mu_{m,n}$
is a star-homomorphism, the positive version uses self-adjoint projections
and matrix adjoint as its dagger.  Its comparison with bare right modules is
algebraic: the bare module category carries no chosen Hilbert form or dagger.
\end{definition}
\begin{scope}
The positive construction requires star-preserving structural maps.  An
algebraic module equivalence does not manufacture a dagger on its target.
\end{scope}
\provenance{D1223}{theory/sidequests/f1-limit/universal.md}{finite matrix-corner derivation}{theory/verdicts/f1-limit-adjudication.md}

\begin{definition}[Positive Hecke sequence and its universal completion]
For real $q>0$, put $H_0(q)=H_1(q)=\mathbb C$.  For $n\geq2$, let $H_n(q)$
be generated by self-adjoint $T_1,\ldots,T_{n-1}$ with
\[
 (T_i-q)(T_i+1)=0,
 \quad T_iT_{i+1}T_i=T_{i+1}T_iT_{i+1},
 \quad T_iT_j=T_jT_i\quad(|i-j|>1).
\]
For $w\in S_n$, $T_w$ is the reduced-word basis element and
$\tau_n(T_w)=\delta_{w,1}$.  Ordered block permutations define
\[
 \iota_{m,n}(T_u\otimes T_v)=T_{u\times v}.
\]
These maps are injective unital trace-preserving star-homomorphisms and make
$H_*(q)$ a positive multiplicative sequence.

Let
\[
 R=\mathbb Z[t,t^{-1},\{P_\alpha(t)^{-1}\}_\alpha]
\]
be the localized coefficient ring used by the generic based Hecke algebras.
The generic category is the finite additive self-adjoint-idempotent envelope
\[
 \mathcal U_R=\operatorname{Kar}_*\!\bigl(\operatorname{Add}
                 (\mathbb C[H_*(R)])\bigr).
\]
Evaluation is defined for its displayed objects and arrows.  The full
positive fibre is
\[
 \mathcal U_q=\operatorname{Proj}_*(H_*(q)).
\]
For $x=[1]$, $\operatorname{End}_{\mathcal U_q}(x^{\otimes n})=H_n(q)$.
A nonzero projection corner uses the coefficient trace divided by the trace
of its corner unit.  Normalized states and effects live on these finite
$C^*$-endomorphism algebras.
\end{definition}
\begin{scope}
The generic envelope contains self-adjoint projectors presented over $R$.
It need not evaluate essentially surjectively onto projections constructed
separately by functional calculus in one positive fibre.  This is an
amplitude category; its operational completion is supplied separately and
does not require a fusion hypothesis.
\end{scope}
\provenance{D1101,D1102,D1224,F1-HCK-POS,F1-HCK-TOWER}{theory/sidequests/f1-limit/universal.md}{theory/checks/f1\_operational\_check.py (H13, finite corner samples)}{theory/verdicts/f1-limit-adjudication.md}

\begin{definition}[Partial-flag completion]
For a composition $\alpha=(\alpha_1,\ldots,\alpha_r)$ of $n$, let
$W_\alpha\leq S_n$ preserve its consecutive blocks and define
\[
 P_\alpha(q)=\sum_{w\in W_\alpha}q^{\ell(w)},\qquad
 e_\alpha(q)=P_\alpha(q)^{-1}\sum_{w\in W_\alpha}T_w.
\]
The partial-flag category has
\[
 \operatorname{Hom}(\alpha,\beta)=e_\beta H_n(q)e_\alpha
\]
for equal totals, zero Homs otherwise, identity $e_\alpha$, adjoint inherited
from $H_n(q)$, and concatenation tensor.  Its projection-model functor sends
\[
 \alpha\longmapsto e_\alpha H_n(q),\qquad
 z\longmapsto L_z:e_\alpha H_n(q)\to e_\beta H_n(q).
\]
The extension under finite sums and self-adjoint retracts is the universal
positive Hecke category above.
\end{definition}
\begin{scope}
The normalization by $P_\alpha$ is essential.  The completion statement is
proved below from the complete-flag object, not assumed from an older sketch
about partial flags.
\end{scope}
\provenance{D1141,D1225}{theory/sidequests/f1-limit/universal.md}{theory/checks/f1\_operational\_check.py (H13, levels 3 and 4 at the declared parameters)}{theory/verdicts/f1-limit-adjudication.md}

\begin{definition}[Davydov--Molev parameter and finite module category]
Put $v=\sqrt q>0$.  If $t_i^{\mathrm{DM}}$ satisfies
\[
 (t_i^{\mathrm{DM}}-v)(t_i^{\mathrm{DM}}+v^{-1})=0,
\]
then the project conventions are
\[
 q_{\mathrm{project}}=v^2,
 \qquad T_i^{\mathrm{project}}=v,t_i^{\mathrm{DM}}.
\]
Let $\operatorname{DM}_{\mathrm{fin}}(v)$ be the finite-support,
finite-dimensional right-module subcategory of the associated Hecke
Schur--Weyl category.  Its tensor is the induced module just defined.  The
projection realization $(r,e)\mapsto eH_n(q)^r$ is the proposed strong
monoidal comparison
\[
 \mathcal U_q\simeq\operatorname{DM}_{\mathrm{fin}}(\sqrt q).
\]
\end{definition}
\begin{scope}
This comparison is algebraic.  It respects the multiplicative sequence and
induction tensor, but it asserts no dagger on the bare-module target.
\end{scope}
\provenance{D1226,D1227}{theory/sidequests/f1-limit/universal.md}{parameter identity and structured proof; no dedicated exact gate}{theory/verdicts/f1-limit-adjudication.md}

\begin{theorem}[Universal finite Hecke composition]
{\normalfont\statusProved}\quad For every $q>0$, the category $\mathcal U_q$
is algebraically strong-monoidally equivalent to the finite-support category
of finite-dimensional right modules over the Hecke multiplicative sequence,
with tensor
\[
 (M\otimes N)\otimes_{H_m(q)\otimes H_n(q)}H_{m+n}(q).
\]
It is also the additive self-adjoint Karoubi completion of the partial-flag
corner category.  Generic evaluation applies to the self-adjoint projectors
presented over $R$.
\end{theorem}
\begin{scope}
The bare-module equivalence is algebraic.  Generic evaluation is not claimed
essentially surjective on all fibrewise functional-calculus projections, and
no cross-degree amplitude morphisms are added.
\end{scope}
\provenance{F1-LIM-UNIV,D1221,D1222,D1223,D1224,D1225}{theory/sidequests/f1-limit/universal.md}{theory/checks/f1\_operational\_check.py (H13, finite parabolic corners only)}{theory/verdicts/f1-limit-adjudication.md}
\begin{proof}
Contravariance gives $(m\cdot a)\cdot b=m\cdot(ab)$, fixing right-module
variance.  Both bracketings of three induced modules reduce canonically to
\[
 (L\otimes M\otimes N)
 \otimes_{A_l\otimes A_m\otimes A_n}A_{l+m+n},
\]
so associativity and the pentagon follow from the common balanced universal
map.  The representables multiply because
$((A_m\otimes A_n)\otimes_{A_m\otimes A_n}A_{m+n})\cong A_{m+n}$.

Finite sums of representables give free right modules; splitting an
idempotent $e$ gives $eA_n^r$, and
$\operatorname{Hom}_{A_n}(eA_n^r,fA_n^s)=fM_{s,r}(A_n)e$.  Since a
finite-dimensional $C^*$-algebra is semisimple, all of its finite-dimensional
modules are projective.  This proves the algebraic equivalence.  Star
preservation of the block maps sends tensor projections to projections.

For partial flags, write $x_\alpha=\sum_{w\in W_\alpha}T_w$.  Pairing terms
along a simple parabolic reflection gives $T_sx_\alpha=qx_\alpha$ and its
right-handed star version.  Hence
\[
 x_\alpha^2=P_\alpha x_\alpha,\qquad x_\alpha^*=x_\alpha,\qquad
 \tau_n(e_\alpha)=P_\alpha^{-1}.
\]
Block length additivity gives
$P_{\alpha\mathbin{\|}\beta}=P_\alpha P_\beta$ and
$\iota(e_\alpha\otimes e_\beta)=e_{\alpha\mathbin{\|}\beta}$.
Finally the complete composition $(1,\ldots,1)$ has $e=1$ in each degree,
so its finite sums and retracts generate every object of $\mathcal U_q$.
\end{proof}

\begin{theorem}[Algebraic Hecke Schur--Weyl comparison]
{\normalfont\statusProved}\quad With
$q_{\mathrm{project}}=v^2$ and
$T_i^{\mathrm{project}}=v,t_i^{\mathrm{DM}}$, the category $\mathcal U_q$
is algebraically strong-monoidally equivalent to
$\operatorname{DM}_{\mathrm{fin}}(v)$.
\end{theorem}
\begin{scope}
Only finite degree support and finite-dimensional modules are included.  The
comparison does not turn the algebraic Hecke braiding into a physical unitary
or endow the bare-module category with a dagger.
\end{scope}
\provenance{F1-LIM-DM,D1226,D1227}{theory/sidequests/f1-limit/universal.md}{exact parameter calculation; categorical equivalence by proof}{theory/verdicts/f1-limit-adjudication.md}
\begin{proof}
Substitution $t_i^{\mathrm{DM}}=v^{-1}T_i$ turns its quadratic relation into
$(T_i-v^2)(T_i+1)=0$.  The braid and distant-commutation relations are
homogeneous, and the same scaling respects every ordered block inclusion.
Thus the full multiplicative sequences agree.  The preceding theorem then
identifies their finite right-module categories with their induction tensor.
\end{proof}

\subsection{Algebraic braids and unitary cactus exchange}

\begin{definition}[Adjacent spectral sign]
In $H_n(q)$ define
\[
 u_i=\frac{2T_i+1-q}{q+1}
     =2\frac{T_i+1}{q+1}-1.
\]
This is $+1$ on the $q$-eigenspace of $T_i$ and $-1$ on its
$(-1)$-eigenspace.  Its adjacent braid defect is
\[
 u_i u_{i+1}u_i-u_{i+1}u_i u_{i+1}
 =-\frac{(q-1)^2}{(q+1)^2}(u_i-u_{i+1}).
\]
\end{definition}
\begin{scope}
The definition is fibrewise for $q>0$.  It produces local unitary exchanges,
not a braiding away from one.
\end{scope}
\provenance{D1228}{theory/sidequests/f1-limit/exchange.md}{theory/checks/f1\_limit\_check.py (L7, six declared rational parameters)}{theory/verdicts/f1-limit-adjudication.md}

\begin{definition}[Polar longest-element reversal]
Let $w_0^{(n)}$ be the longest permutation and $D_n=T_{w_0^{(n)}}$.  In the
positive fibre put
\[
 |D_n|=(D_n^2)^{1/2},\qquad J_n=D_n|D_n|^{-1},\qquad J_0=J_1=1.
\]
The square root is the unique positive $C^*$-functional-calculus root.  For
an interval $I$ of consecutive tensor positions, $J_I$ is the shifted copy
of $J_{|I|}$.
\end{definition}
\begin{scope}
This definition uses positivity and functional calculus at real $q>0$; it is
not an elementwise construction over the localized generic coefficient ring.
\end{scope}
\provenance{D1229}{theory/sidequests/f1-limit/exchange.md}{theory/checks/f1\_limit\_check.py (L13, ranks at most 4 and q=1,4,9)}{theory/verdicts/f1-limit-adjudication.md}

\begin{definition}[Unitary block commutor]
For $m,n\geq0$, define on tensor powers
\[
 \sigma_{m,n}=J_{m+n}\,\iota_{m,n}(J_m\otimes J_n).
\]
On projection-completed objects, restrict this operator to the source and
target corners; on finite sums also apply the canonical matrix-index flip.
Its coherence consists of naturality, symmetry
$\sigma_{n,m}\sigma_{m,n}=1$, the unit laws, and the cactus square formed by
the interval reversals.
\end{definition}
\begin{scope}
The commutor is ordered and coboundary.  Symmetry here means involutivity of
the commutor, not the braid relation for three adjacent factors.
\end{scope}
\provenance{D1230}{theory/sidequests/f1-limit/exchange.md}{theory/checks/f1\_limit\_check.py (L13, finite cactus samples)}{theory/verdicts/f1-limit-adjudication.md}

\begin{definition}[Analytic domain of physical exchange]
The block commutor is extra canonical structure on every positive real
$C^*$-fibre and varies continuously for $q>0$.  It is the polar unitary of
the positive minimal Hecke braid exchanging the two blocks.  At $q=1$ it is
the ordinary unitary block permutation.  Its functional-calculus
coefficients need not belong to the localized generic ring.
\end{definition}
\begin{scope}
No algebraic generic commutor is stipulated, and no braid coherence is
claimed for $q\ne1$.
\end{scope}
\provenance{D1231}{theory/sidequests/f1-limit/exchange.md}{theory/checks/f1\_limit\_check.py (L13, declared fibres only)}{theory/verdicts/f1-limit-adjudication.md}

\begin{theorem}[Spectral unitaries do not braid away from one]
{\normalfont\statusProved}\quad For $q>0$, the standard algebraic Hecke
exchange is invertible but unitary for the coefficient-trace $C^*$-structure
only at $q=1$.  The adjacent spectral signs are self-adjoint unitaries with
the exact defect displayed above, and therefore do not braid when $q\ne1$.
\end{theorem}
\begin{scope}
This separates algebraic Yang--Baxter exchange from physical unitary
exchange.  It does not obstruct a coboundary commutor.
\end{scope}
\provenance{F1-LIM-EXCH,D1226,D1228}{theory/sidequests/f1-limit/exchange.md}{theory/checks/f1\_limit\_check.py (L7, six rational parameters)}{theory/verdicts/f1-limit-adjudication.md}
\begin{proof}
The algebraic generator $v^{-1}T_i$ has eigenvalues $v$ and $-v^{-1}$,
so both have modulus one exactly at $v=1$.  Meanwhile
$e_i=(T_i+1)/(q+1)$ is a self-adjoint projection and $u_i=2e_i-1$.
Writing $u_1=ax+b$, $u_2=ay+b$, with
$a=2/(q+1)$, $b=(1-q)/(q+1)$ and $x=T_1,y=T_2$, braid cancellation gives
\[
 u_1u_2u_1-u_2u_1u_2=a^2b(x^2-y^2)+ab^2(x-y).
\]
Since $x^2-y^2=(q-1)(x-y)$, this is the stated multiple of $u_1-u_2$.
The standard basis makes $T_1-T_2$ nonzero.
\end{proof}

\begin{theorem}[Canonical unitary cactus exchange]
{\normalfont\statusProved}\quad For every $q>0$, the polar reversals $J_n$
define a continuous natural unitary coboundary commutor on all of
$\mathcal U_q$.  It is symmetric, satisfies the cactus relations, restricts
to every partial-flag and arbitrary projection corner, is the polar unitary
of the block braid, and becomes ordinary symmetric exchange at $q=1$.
\end{theorem}
\begin{scope}
The theorem is analytic on positive fibres and positive intervals.  It makes
no claim over the generic localized ring and no claim of braiding away from
one.
\end{scope}
\provenance{F1-LIM-COB,D1229,D1230,D1231}{theory/sidequests/f1-limit/exchange.md}{theory/checks/f1\_limit\_check.py (L13, n at most 4 and q=1,4,9)}{theory/verdicts/f1-limit-adjudication.md}
\begin{proof}
The longest element is self-adjoint and invertible.  The relation
$w_0s_i=s_{n-i}w_0$ gives
$D_nT_i=T_{n-i}D_n$, hence $D_n^2$ is central.  Its positive modulus is
therefore central, so $J_n$ is a canonical self-adjoint unitary satisfying
$J_nT_iJ_n=T_{n-i}$.  Shifted copies on disjoint intervals commute, and
unitary conjugation carries a subinterval reversal to its reflected
subinterval reversal.  Thus
\[
 J_{[p,r]}J_{[k,l]}
 =J_{[p+r-l,p+r-k]}J_{[p,r]},
\]
which is the interval presentation of the cactus relations.

The product $J_{m+n}(J_m\otimes J_n)$ reverses internally and then totally,
so the internal reversals cancel and the two ordered block algebras swap.
This proves naturality; the same interval identities prove symmetry and the
cactus square.  Conjugation carries any tensor projection to its swapped
projection, so compression extends the commutor to all finite sums and
retracts.

If $B_{m,n}=D_{m+n}(D_m\otimes D_n)^{-1}$ is the positive minimal block
braid, then
\[
 \sigma_{m,n}^*B_{m,n}=|D_{m+n}|\,|D_m\otimes D_n|^{-1}>0.
\]
Uniqueness of polar decomposition identifies $\sigma_{m,n}$ as its polar
unitary.  Functional calculus varies continuously on positive invertible
matrices.  At one, every $T_i$ is a transposition, so the construction is the
ordinary block permutation.
\end{proof}
